\sectionTitle{经历}{\faCode}

\begin{experiences}
\experience
   {2023年12月}   {安徽大学国际脑科学工程研究中心}{研究实习}{2025年03月}
    {\textbf{项目导师:} \link{https://xingbod.github.io}{董兴波老师}。\\ \textbf{贡献:} 开发超声、PET 与皮肤影像的重建和泛化方法；参与多篇已发表或录用论文，并担任 DeepPET 项目负责人，获国家级和省级创新创业奖项。}
    {}

\experience
   {2025年07月}   {阿尔伯塔大学 Vision and Learning Lab}{研究实习}{2025年10月}
  {\textbf{项目导师:} \link{http://www.ece.ualberta.ca/~lcheng5/}{Li Cheng教授}。\\ \textbf{贡献:} 开发融合物理先验与深度学习的 3D CT 重建方法，面向低剂量、有限角度和高噪声场景，提升受限采集条件下的图像质量。}
  {}

\experience
   {2025年12月}   {上海交通大学转化医学研究院}{研究实习}{2026年10月}
  {\textbf{项目导师:} \link{https://lei.bi/}{Lei Bi教授}。\\ \textbf{贡献:} 开展基于视觉语言模型的多模态临床智能和三维医学重建工作，支持医学人工智能的临床转化。}
  {}

\end{experiences}



  %   {2024年5月}   {面向异质超声数据的智能诊断模型研究}{项目负责人}
  %   {
  %                     \begin{itemize}
  %                       \item \textbf{项目描述:} 创建了一个大规模、公开可访问的超声数据集，用于改进机器学习模型在超声图像中的疾病分类能力。超声广泛用于内部结构的诊断，但操作员差异、噪声和视野受限等因素使得诊断一致性存在挑战。该数据集包含来自多种来源的1833张匿名超声图像，涵盖13种异常，支持疾病分类模型在多种环境下的基准测试，助力AI辅助超声诊断的研究进展。
  %                       \item \textbf{关键贡献:} 作为\textbf{第一作者}撰写了论文《An Ultrasound Dataset in the Wild for Machine Learning of Disease Classification》，拟投稿至\textbf{《Scientific Data》 (修订中) (Springer Nature, JCR Q1)}。
  %                       \item \faGithub: 
  %                       \link{https://github.com/Bean-Young/US-Dataset}{项目Github链接}。              
  %                     \end{itemize}
  %                   }
  %                   {机器学习,超声数据集,异质性,公开数据}
  % \emptySeparator
  % \experience
  % {2024年1月}   {基于学习的脑PET图像重建方法及可解释性研究}{项目负责人}
  %   {
  %                     \begin{itemize}
  %                       \item \textbf{项目描述:} 开发了DeepPET，一种结合物理建模与深度学习的创新PET图像重建方法，可在显著降低辐射剂量的同时提高图像质量。该方法提高了重建速度，减少了传感器需求，从而降低了成本。通过生成模型，DeepPET将低剂量图像转换为高剂量、无伪影图像，推进PET成像的效率，助力早期诊断和可靠临床使用。
  %                       \item \textbf{关键贡献:} 作为\textbf{项目负责人}，通过中国大学生创新创业项目获得\textbf{国家级立项}。带领团队作为\textbf{项目负责人}在中国国际“互联网+”大学生创新创业大赛安徽省赛中荣获\textbf{金奖}。作为\textbf{核心开发者}参与申请了一项\textbf{发明专利}《基于先验图像的PET图像重建方法及三维感知方法》。
  %                       \item \faGithub: \link{https://github.com/Bean-Young/PET-Project}{项目Github链接}。              
  %                     \end{itemize}
  %                   }
  %                   {正电子发射断层扫描,低成本,图像重建,降噪}
  % \emptySeparator
  % \experience
  % {2024年8月}   {基于3D Gaussian Splatting的超声图像渲染方法}{项目负责人}
  %   {
  %                     \begin{itemize}
  %                       \item \textbf{项目描述:} 基于3D Gaussian Splatting开发了一种高保真超声图像渲染方法。该方法实现了超声图像的快速渲染，并在多个数据集上达到了最先进的性能。
  %                       \item \textbf{关键贡献:} 作为\textbf{第一作者}撰写了一篇关于该方法的论文，预计于12月初步完成。           
  %                     \end{itemize}
  %                   }
  %                   {三维视觉,超声图像处理,视角生成,体素渲染}


  %     \textbf{领域泛化研究:}
%     \begin{itemize}
%         \item 《EM-Net: Effective and Morphology-aware Network for Skin Lesion Segmentation》，发表在\textbf{《Expert Systems with Applications》(Elsevier Ltd, JCR Q1)}。提出基于形态学感知的网络结构，有效解决皮肤损伤分割中的领域偏移问题。Link: \href{https://github.com/Bean-Young/EM-Net}{\faGithub} \href{https://arxiv.org/abs/2401.xxxxx}{\faFilePdfO} \href{https://your-project-website.com}{\faGlobe}
        
%         \item 《Abnormality-aware Prompting for Test-Time Adaptation in Breast Ultrasound Segmentation》(已投NeurIPS)，针对乳腺超声图像分割中的域适应挑战，设计异常感知提示机制，提高模型在未见域上的泛化能力。
        
%         \item 《Federated Learning-Based Virtual Dual-Energy CT Generation from Single-Energy CT for Gout Detection》(已投Digital Health)，创新性地将联邦学习应用于虚拟双能CT生成，保护患者隐私的同时提高痛风检测准确率。
%     \end{itemize}
    
%     \textbf{医学图像重建:}
%     \begin{itemize}
%         \item 《Depth-Aware Gaussian Splatting with Propagation Properties for Ultrasound Rendering》(已投BMVC)，首次将高斯散射技术应用于超声图像三维重建，实现深度感知的超声图像渲染，提升三维可视化质量。
        
%         \item 《Explicit and Implicit Representations in AI-based 3D Reconstruction for Radiology: A Systematic Review》(已投Medical Image Analysis)，系统梳理医学影像三维重建领域的表示学习方法，为未来研究提供理论指导。
%     \end{itemize}
    
%     \textbf{超声成像研究:}
%     \begin{itemize}
%         \item 《An annotated heterogeneous ultrasound database》，发表在\textbf{《Scientific Data》(Nature旗下期刊，JCR Q1)}，构建首个大规模异构超声数据库，包含多器官、多设备的标注数据，为超声AI研究提供基础支撑。
        
%         \item 《Auto-US: An Agent for Ultrasound Video Diagnosis Using Video Classification Framework and LLMs》(已投npj artificial intelligence)，开发基于大模型的超声视频诊断智能体系统，融合视频分类框架与大语言模型，实现超声诊断的智能辅助。
%     \end{itemize}
    
%     \textbf{专利及项目成果:}
%     \begin{itemize}
%         \item 《基于先验图像的PET图像重建方法及PET图像3D感知方法》，国家发明专利，CN118411435A，实质审查中。
        
%         \item 《低剂量PET图像处理的3D高斯模型生成和渲染方法》，国家发明专利，CN119625190A，实质审查中。

%         \item 《DeepPET:低放射高成像质量的智能脑PET设备》，中国国际大学生创新大赛， 安徽省金奖。

%         \item 《基于学习的脑PET图像重建方法及可解释性研究》，大学生创新创业训练计划，国家级优秀结题，安徽大学“创新之星”。
%     \end{itemize}


    % \item \textbf{研究经历:} 在这段研究经历中，我专注于领域泛化问题和医学图像重建研究，涉及超声、PET和皮肤损伤等多种医学成像模态。致力于解决多源异构医学数据的处理与分析难题，推动医学影像AI技术的临床实用化。
    
    % \item \textbf{关键贡献:} 作为\textbf{第一/共一作者}撰写了7篇高水平科研论文，并申请2项国家发明专利。同时，作为学生项目负责人带领团队获得国家级大学生创新创业训练计划优秀结题、中国国际大学生创新大赛安徽省金奖，并荣获安徽大学"创新之星"称号。
